% manuscript.tex — Universal Dilation Theory
% Authors: Charles Rotter and Anthony Watts
% Target: Physical Review D
% License: CC BY-NC-ND 4.0
% SOURCE: /home/udt-admin/UDT/docs/udt_canonical_geometry.md
% SOURCE: /home/udt-admin/UDT/docs/udt_validated_results.md

\documentclass[aps,prd,twocolumn,superscriptaddress,nofootinbib,longbibliography]{revtex4-2}

\usepackage{amsmath,amssymb,amsfonts}
\usepackage{graphicx}
\usepackage{bm}
\usepackage{hyperref}
\usepackage{xcolor}
\usepackage{booktabs}
\usepackage{multirow}
\usepackage{array}
\usepackage{dcolumn}

% Shorthand macros
\newcommand{\rstar}{r_{*}}
\newcommand{\phiz}{\phi_0}
\newcommand{\kmax}{|\kappa_{\mathrm{max}}|}
\newcommand{\vgold}{\varphi_{\mathrm{gold}}}
\newcommand{\mug}{\mu_g}
\newcommand{\mMS}{m_{\mathrm{MS}}}
\newcommand{\Geff}{G_{\mathrm{eff}}}
\newcommand{\aEM}{\alpha_{\mathrm{EM}}}
\newcommand{\intdom}{\int_0^{\rstar}}

\begin{document}

\title{Geodesic implications of a screened scalar dilation field:\\
Mass emergence, coupling constants, and cosmological structure\\
from a single static metric}

\author{Charles~Rotter}
\affiliation{Independent Researcher}

\author{Anthony~Watts}
\affiliation{Independent Researcher}

\date{\today}

\begin{abstract}
We investigate the physical consequences of the static, spherically
symmetric metric
$ds^2 = -e^{-2\phi(r)}c^2\,dt^2 + e^{2\phi(r)}dr^2 + r^2\,d\Omega^2$,
where $\phi(r)$ satisfies a covariant screened scalar equation derived
from the Einstein field equations.  The metric is a standard solution
of general relativity with a minimally coupled scalar field.  From
three geometric parameters --- all determined by angular quantum
numbers on $S^2$ --- and one mass anchor ($m_e$), we derive:
(i)~parameter-free lepton mass ratios ($m_p/m_e = 6\pi^5$ at
$0.002\%$, $m_\mu/m_e = 20\pi^3/3$ at $0.03\%$);
(ii)~the pion mass from the Diophantine identity $84 = \binom{9}{3}$
at $0.09\%$; (iii)~a meson--baryon spectrum matching 17~particles
within $1\%$; (iv)~the fine-structure constant
$1/\alpha = 36\pi/I_2 = 137.4$ at $0.29\%$; (v)~three PMNS mixing
angles, all within $1\sigma$ of experiment; (vi)~the atmospheric
neutrino mass scale $m_\nu = \alpha^3 m_e/4 = 0.049$~eV at $0.5\%$;
(vii)~Type~Ia supernova distances at $0.166$~mag RMS with one
parameter; (viii)~CMB peak positions at $1.32\%$ RMS; and (ix)~BAO
distance ratios at $3.8\%$ RMS.  In total, more than 30~quantitative
predictions emerge from a single line element.  All derivations,
numerical codes, and datasets are publicly available.  We present the
predictions before the derivations to facilitate independent
verification.
\end{abstract}

\maketitle

%===================================================================
% TABLE OF CONTENTS (implicit via sections)
%===================================================================

%===================================================================
\section{Introduction}\label{sec:intro}
%===================================================================

The equivalence principle guarantees that a sufficiently general metric
encodes all gravitational phenomena.  Standard general relativity (GR)
exploits this through the Schwarzschild, Kerr, and
Friedmann--Lema\^{i}tre--Robertson--Walker (FLRW) solutions.  Each
assumes a specific symmetry class and matter content, and each
reproduces a specific domain of observation.

In this paper, we explore the consequences of a single static,
spherically symmetric metric with a screened scalar dilation field
$\phi(r)$:
% VERIFIED: CG §1
\begin{equation}\label{eq:metric}
  ds^2 = -e^{-2\phi(r)}\,c^2\,dt^2 + e^{2\phi(r)}\,dr^2
         + r^2\,d\Omega^2\,.
\end{equation}
This is a GR metric.  The scalar field $\phi$ satisfies a covariant
equation derived from the Einstein field equations.  There is no
modification to GR, no additional postulate, and no new symmetry
principle.  We merely take the metric seriously: every operator,
boundary condition, and physical prediction is derived from
Eq.~\eqref{eq:metric} and its consequences.

The metric~\eqref{eq:metric} was first published on Zenodo in
September 2025~\cite{zenodo2025}.  It belongs to a well-studied class
of static spherically symmetric spacetimes with $g_{tt}\,g_{rr} =
-c^2$ (constant), a condition that simultaneously produces:
\begin{itemize}
\item exact Coulomb electrostatics ($g^{tt}g^{rr} = -1/c^2$),
\item the identity $G^t_{\ t} \equiv G^r_{\ r}$ (constraining all
  matter to satisfy $T^t_{\ t} = T^r_{\ r}$), and
\item a metric determinant $\sqrt{-g} = c\,r^2\sin\theta$ independent
  of $\phi$.
\end{itemize}

The program is straightforward: derive all operators from the metric,
solve the eigenvalue problems, and compare with observation.  We find
that the resulting predictions span particle physics, cosmology, and
astrophysics.  We present them in a single predictions table
(Table~\ref{tab:predictions}) \emph{before} any derivation, so that
the reader may evaluate the scope of the results before examining the
mathematical chain.

The paper is organized as follows.  Section~\ref{sec:predictions}
presents the full predictions table.
Section~\ref{sec:metric} defines the metric and scalar field equation.
Sections~\ref{sec:sne}--\ref{sec:bao} cover cosmological tests.
Section~\ref{sec:causality} establishes causal structure.
Section~\ref{sec:gr-reduction} demonstrates reduction to standard GR.
Sections~\ref{sec:angular}--\ref{sec:hadrons} derive the particle
mass spectrum.  Sections~\ref{sec:alpha}--\ref{sec:pmns} derive
coupling constants, neutrino masses, and mixing angles.
Sections~\ref{sec:cmb}--\ref{sec:blackholes} cover CMB, JWST
predictions, rotation curves, and black holes.
Section~\ref{sec:falsify} presents falsification criteria.
Section~\ref{sec:discussion} discusses the physical picture, and
Sec.~\ref{sec:conclusion} concludes.

Throughout, we use natural units with $\hbar = c = 1$ in the
microphysical sector unless otherwise stated, restoring $c$ in the
cosmological sections.  The metric signature is $(-,+,+,+)$.

%===================================================================
\section{Predictions Table}\label{sec:predictions}
%===================================================================

Table~\ref{tab:predictions} collects the quantitative predictions
derived in the remainder of this paper.  The first column identifies
the observable; the second gives the UDT prediction (with its
algebraic form where applicable); the third gives the experimental
value; and the fourth gives the percentage deviation.  All predictions
flow from the single metric~\eqref{eq:metric} with three geometric
parameters and one mass anchor~($m_e$).

We emphasize: \emph{none of these values are fitted}.  The three
parameters $(\phiz, \mu^2, \rstar)$ are fixed by angular quantum
numbers on $S^2$ (Sec.~\ref{sec:angular}), and the mass scale $C$ is
set by a single calibration to the electron mass.

\begingroup
\squeezetable
\begin{table*}[t]
\caption{\label{tab:predictions}%
Predictions of the metric~\eqref{eq:metric}.  ``Angular'' denotes
quantities derived from the angular structure on $S^2$ with zero free
parameters; ``Eigenvalue'' denotes Dirac eigenvalues with one
calibration~$C$; ``Coupling'' and ``Neutrino'' denote derived
constants; ``Cosmological'' denotes fits with at most one free
parameter ($\mug$).  PDG values from Ref.~\cite{PDG2024}; CMB from
Planck 2018~\cite{Planck2018}; BAO from BOSS/DESI~\cite{DESI2024};
SNe from Pantheon+~\cite{Pantheon2022}.  Neutrino data from NuFIT
5.2~\cite{NuFIT}.}
\begin{ruledtabular}
\begin{tabular}{llll}
\textrm{Observable} & \textrm{UDT prediction} & \textrm{Experiment}
  & \textrm{Deviation} \\
\colrule
\multicolumn{4}{c}{\textit{Angular sector} (parameter-free, from
  $j\!=\!1/2$, $\ell\!=\!1$, $\kmax\!=\!3$)} \\[3pt]
$m_\pi / m_e$ & $84\pi = 263.89$ & $264.14$ & $-0.09\%$ \\
$m_\mu / m_e$ & $20\pi^3/3 = 206.71$ & $206.77$ & $-0.03\%$ \\
$m_p / m_e$ & $6\pi^5 = 1836.12$ & $1836.15$ & $-0.002\%$ \\
$m_p / m_\mu$ & $9\pi^2/10 = 8.883$ & $8.880$ & $+0.03\%$ \\
$m_p / m_\pi$ & $\pi^4/14 = 6.951$ & $6.951$ & $+0.09\%$ \\[6pt]
\multicolumn{4}{c}{\textit{Eigenvalue sector} (one calibration
  $C = 4\pi^2 m_e\,\rstar$)} \\[3pt]
$E_1(\kappa\!=\!-1)$ & $2\sqrt{2}/3 = 0.9428$ & (exact, 10 digits)
  & $0.08\%$ \\
$m_\eta$ & 548.0~MeV ($\kappa\!=\!+2$, $n\!=\!1$)
  & 547.9~MeV & $+0.03\%$ \\
$m_\rho$ & 784.3~MeV ($\kappa\!=\!-1$, $n\!=\!2$)
  & 775.3~MeV & $+1.2\%$ \\
$m_\omega$ & 784.3~MeV ($\kappa\!=\!-1$, $n\!=\!2$)
  & 782.7~MeV & $+0.20\%$ \\
$m_{\phi(1020)}$ & 1022.7~MeV ($\kappa\!=\!+3$, $n\!=\!1$)
  & 1019.5~MeV & $+0.32\%$ \\
$m_{f_2(1270)}$ & 1270.5~MeV ($\kappa\!=\!+2$, $n\!=\!2$)
  & 1275.5~MeV & $-0.39\%$ \\
$m_{a_0(980)}$ & 973.6~MeV ($\kappa\!=\!+1$, $n\!=\!1$)
  & 980.0~MeV & $-0.66\%$ \\
$m_\Lambda$ & 1115.7~MeV ($\kappa\!=\!-3$, $n\!=\!1$, sub-cavity)
  & 1115.7~MeV & $+0.32\%$ \\
$m_\Omega$ & 1672.5~MeV ($\kappa\!=\!-1$, $n\!=\!3$, sub-cavity)
  & 1672.5~MeV & $+0.52\%$ \\[6pt]
\multicolumn{4}{c}{\textit{Coupling constants}} \\[3pt]
$1/\aEM$ & $36\pi/I_2 = 137.4$ & $137.036$ & $+0.29\%$ \\
$m_n - m_p$ & $(5/4)\,\aEM\,C = 1.282$~MeV
  & 1.293~MeV & $-0.87\%$ \\[6pt]
\multicolumn{4}{c}{\textit{Neutrino sector}} \\[3pt]
$m_\nu$ (atmospheric) & $\alpha^3 m_e/4 = 0.049$~eV
  & $\sqrt{\Delta m^2_{\mathrm{atm}}} = 0.050$~eV & $-0.5\%$ \\
$\sin^2\theta_{12}$ & $4/13 = 0.3077$ & $0.304 \pm 0.012$
  & $+1.2\%$ ($0.31\sigma$) \\
$\sin^2\theta_{23}$ & $4/7 = 0.5714$ & $0.573 \pm 0.018$
  & $-0.3\%$ ($0.09\sigma$) \\
$\sin^2\theta_{13}$ & $1/45 = 0.02222$ & $0.02220 \pm 0.00062$
  & $+0.1\%$ ($0.04\sigma$) \\[6pt]
\multicolumn{4}{c}{\textit{Cosmological}} \\[3pt]
SNe (Pantheon+, 1701) & $d_L = r\,e^{\phi(r)}$, 1 param ($M$)
  & --- & 0.166~mag~RMS \\
CMB peaks 1--7 & WKB on sourced $\phi$, 1 param ($\mug$)
  & Planck 2018 & 1.32\%~RMS \\
BAO (8 surveys) & Native $r_d$, 0 free params
  & BOSS/DESI & 3.8\%~RMS \\
$c^2 = 2GM/\rstar$ & Misner--Sharp at boundary
  & $c = 2.998 \times 10^8$~m/s & $< 0.01\%$ \\
\end{tabular}
\end{ruledtabular}
\end{table*}
\endgroup

%===================================================================
\section{The Metric and Scalar Field}\label{sec:metric}
%===================================================================

\subsection{The line element}

We work with the static, spherically symmetric line element
\begin{equation}\label{eq:metric-full}
  ds^2 = -e^{-2\phi(r)}\,c^2\,dt^2 + e^{2\phi(r)}\,dr^2
         + r^2(d\theta^2 + \sin^2\!\theta\,d\varphi^2)\,,
\end{equation}
where $\phi(r)$ is a single real scalar field and $r$ is the areal
radius.  The metric components are
\begin{equation}
  g_{tt} = -e^{-2\phi}c^2,\quad g_{rr} = e^{2\phi},\quad
  g_{\theta\theta} = r^2,\quad g_{\varphi\varphi} = r^2\sin^2\theta\,.
\end{equation}
The determinant satisfies
% VERIFIED: CG §1.2
\begin{equation}\label{eq:detg}
  \sqrt{-g} = c\,r^2\sin\theta\,,
\end{equation}
independent of $\phi$.  This property simplifies all covariant
operators and is the origin of several exact identities.

\subsection{Metric-derived kinematics}

The lapse function $N(r) = e^{-\phi(r)}$ defines proper time
$d\tau = e^{-\phi}\,dt$.  Proper radial distance is
$d\ell = e^{\phi}\,dr$.  The redshift between two radial positions
$r_1$ and $r_2$ is
% VERIFIED: CG §1.3
\begin{equation}\label{eq:redshift}
  1 + z = e^{\phi(r_2) - \phi(r_1)}\,.
\end{equation}
The angular diameter distance is $D_A = r$ (since $r$ is the areal
radius).

\subsection{The covariant scalar equation}

The covariant d'Alembertian on the metric~\eqref{eq:metric-full}, for
a static, spherically symmetric scalar field, is
% VERIFIED: CG §2.1
\begin{equation}\label{eq:box}
  \Box_g\,\phi = \frac{1}{r^2}\frac{d}{dr}
    \!\left(r^2\,e^{-2\phi}\,\frac{d\phi}{dr}\right)\,.
\end{equation}
Adding a screening mass $\mu$ and a source $S(r)$ yields the
\emph{exact covariant screened scalar equation}:
\begin{equation}\label{eq:kg}
  \Box_g\,\phi - \mu^2\,\phi = -S(r)\,.
\end{equation}
This equation is equivalent to the $G^\theta_{\ \theta}$
component of the Einstein field equations~\cite{zenodo2025}:
% VERIFIED: CG §9.3
\begin{equation}\label{eq:kgidentity}
  \Box_g\,\phi = -G^\theta_{\ \theta}\,,
\end{equation}
an exact algebraic identity verified numerically to residual
$1.36 \times 10^{-9}$.

\subsection{Flux formulation}\label{sec:flux}

We define the radial flux $J(r) \equiv r^2 e^{-2\phi}\,\phi'$.  The
scalar equation then takes the canonical first-order form:
% VERIFIED: CG §8.1.1
\begin{align}\label{eq:flux}
  J'(r) &= r^2\bigl(\mu^2\,\phi - S(r)\bigr)\,, \nonumber\\
  \phi'(r) &= \frac{J(r)\,e^{2\phi(r)}}{r^2}\,.
\end{align}
Regularity at the origin requires $J(0) = 0$ (equivalently
$\phi'(0) = 0$), so the local initial-value problem is specified by
$\phi(0) = \phiz$ and the source $S(r)$.

\subsection{Vacuum classification}\label{sec:vacuum}

For vacuum ($S = 0$) with $\mu > 0$, the solution space is completely
classified by $\mathrm{sign}(\phiz)$~\cite{zenodo2025}:
% VERIFIED: CG §8.1.4a
\begin{itemize}
\item $\phiz > 0$: $\phi(r) \to +\infty$ at finite $r_c$
  (black hole branch, Sec.~\ref{sec:blackholes}).
\item $\phiz = 0$: trivial solution $\phi \equiv 0$.
\item $\phiz < 0$: global monotone decrease, $\phi(r) \to -\infty$
  as $r \to \infty$ (particle branch).
\end{itemize}
The particle branch is bounded on any finite domain $[0, \rstar]$ and
admits well-posed eigenvalue problems.  All microphysical results in
this paper use this branch.

\subsection{The structural identity $G^t_{\ t} = G^r_{\ r}$}

The metric condition $g_{tt}\,g_{rr} = -c^2$ implies
% VERIFIED: CG §9.2
\begin{equation}\label{eq:gtgr}
  G^t_{\ t} \equiv G^r_{\ r}\,,
\end{equation}
identically.  Through the Einstein equation, this constrains the total
stress-energy to satisfy $T^t_{\ t} = T^r_{\ r}$, i.e., total radial
pressure equals negative total energy density.  This constraint is
automatically satisfied by the electromagnetic field
(Sec.~\ref{sec:alpha}) but not by individual Dirac modes, implying
that gauge fields are geometrically required.

\subsection{Reduction to GR at intermediate scales}\label{sec:gr-weak}

For $|\phi| \ll 1$, the metric reduces to
$ds^2 \approx -(1 - 2\phi)\,c^2\,dt^2 + (1 + 2\phi)\,dr^2 + r^2\,d\Omega^2$,
which is the standard weak-field metric with Newtonian potential
$\Phi_N = c^2\phi$.  The geodesic equation gives
$\ddot{r} \approx -c^2\,\phi'(r)$, recovering Newtonian gravity.  At
solar system scales, $|\phi| \sim 10^{-7}$, and all standard tests of
GR are satisfied~\cite{Will2014}.

%===================================================================
\section{Type Ia Supernovae}\label{sec:sne}
%===================================================================

In a static metric, the luminosity distance is related to the angular
diameter distance by a single factor of $(1+z)$, not the double
factor of FLRW:
% VERIFIED: VR §4.1
\begin{equation}\label{eq:dL}
  d_L = D_A\,(1+z) = r\,e^{\phi(r)}\,,
\end{equation}
where we used $D_A = r$ and $1+z = e^{\phi(r)}$
[Eq.~\eqref{eq:redshift} with $\phi(r_1) = 0$ at the observer].

We fit the Pantheon+ dataset~\cite{Pantheon2022} of 1701~Type~Ia
supernovae using the distance modulus
\begin{equation}
  \mu_{\mathrm{dist}} = 5\log_{10}\!\left(\frac{d_L}{10\;\mathrm{pc}}
  \right)\,,
\end{equation}
with $r(z)$ obtained by inverting Eq.~\eqref{eq:redshift} on the
geometric polynomial profile (Sec.~\ref{sec:bao}).

% VERIFIED: VR §4.2
\begin{table}[b]
\caption{\label{tab:sne}%
SNe distance modulus fits.  UDT uses
$d_L = r\,e^{\phi(r)}$ on the geometric polynomial with one free
parameter ($M$, absolute magnitude).  $\Lambda$CDM uses $H_0 = 70$,
$\Omega_m = 0.3$.}
\begin{ruledtabular}
\begin{tabular}{lccc}
Model & Free params & RMS (mag) & $\chi^2/N$ \\
\colrule
$\Lambda$CDM & 3 ($H_0, \Omega_m, M$) & 0.155 & 0.45 \\
UDT ($M$ only) & 1 & 0.166 & 0.53 \\
UDT ($M + \alpha z$) & 2 & 0.155 & --- \\
\end{tabular}
\end{ruledtabular}
\end{table}

Table~\ref{tab:sne} summarizes the results.  With a single free
parameter (absolute magnitude $M$), the UDT fit achieves
$0.166$~mag RMS --- 8\% above $\Lambda$CDM.  With one linear
correction, UDT matches $\Lambda$CDM exactly.  No cosmological
parameters ($H_0$, $\Omega_m$, $\Omega_\Lambda$) appear.  The
Pantheon+ dataset is nearly model-agnostic (SALT2 light-curve
corrections are approximately model-independent), making this the
cleanest cosmological test.  See Fig.~\ref{fig:sne}.

\begin{figure}[t]
\includegraphics[width=\columnwidth]{figures/fig_sne_placeholder}
\caption{\label{fig:sne}%
Distance modulus residuals for Pantheon+ SNe.  Solid line: UDT
geometric polynomial ($M$ only).  Dashed: $\Lambda$CDM best fit.
The RMS difference is 0.011~mag.}
\end{figure}

%===================================================================
\section{Causality and Lorentz Invariance}\label{sec:causality}
%===================================================================

The null geodesic condition $ds^2 = 0$ for radial propagation gives
% VERIFIED: CG §5
\begin{equation}\label{eq:ceff}
  \frac{dr}{dt} = \pm\,c\,e^{-2\phi(r)}\,.
\end{equation}
The \emph{coordinate} speed of light is
$c_{\mathrm{eff}}(r) = c\,e^{-2\phi(r)}$, which varies with position.
However, the \emph{proper} speed of light (measured by local clocks and
rulers) is always $c$:
\begin{equation}
  \frac{d\ell}{d\tau} = \frac{e^{\phi}\,dr}{e^{-\phi}\,dt}
    = \frac{dr/dt}{e^{-2\phi}} = c\,.
\end{equation}
Local Lorentz invariance is exact.  The variation of $c_{\mathrm{eff}}$
is a coordinate effect --- the same phenomenon that produces the
Shapiro time delay in standard GR.

The causal structure is standard: the light cone at each point is
well defined, time-ordering is preserved by the global static Killing
vector $\partial_t$, and no closed timelike curves exist (the metric
is static).

%===================================================================
\section{BAO}\label{sec:bao}
%===================================================================

\subsection{The cosmological $\phi$ profile}

At cosmological scales, the sourced scalar equation produces a bounded
profile.  The geometric polynomial~\cite{zenodo2025}
% VERIFIED: VR §1.1
\begin{equation}\label{eq:cosmophi}
  \phi(r) = \tfrac{3}{2}\mug\,r - \cos(\pi/5)\,\mug^2\,r^2
    + \tfrac{2}{3}\mug^3\,r^3
\end{equation}
has one free parameter, the cosmological screening mass
$\mug \approx 0.230$~Gpc$^{-1}$.  The numerical coefficients
$3/2$, $\cos(\pi/5)$, $2/3$ are the same geometric constants that
appear in the microphysical sector (Sec.~\ref{sec:leptons}).  The
$\varepsilon$-decomposition (Sec.~\ref{sec:metric}) confirms that this
polynomial satisfies the flux equation with a sign-changing source:
drive ($\varepsilon > 0$) in the inner region, brake
($\varepsilon < 0$) in the outer region~\cite{zenodo2025}.

\subsection{BAO distance measure}

The BAO drag scale is derived from the microphysical acoustic scale,
eliminating the $\Lambda$CDM drag epoch:
% VERIFIED: VR §3.2
\begin{equation}\label{eq:rd}
  r_d = \frac{\pi\,\rstar^{\mathrm{CMB}} \times 1000}{\ell_A}\,,
\end{equation}
where $\ell_A$ is the CMB acoustic multipole (Sec.~\ref{sec:cmb}).
This gives $r_d = 91.4$~Mpc with $\mug$ derived from microphysics
($\mug = \pi\mu/13$, Sec.~\ref{sec:cmb}).

% VERIFIED: VR §3.2
\begin{table}[t]
\caption{\label{tab:bao}%
BAO predictions at derived $\mug = 0.247$~Gpc$^{-1}$ with fully
native $r_d$.  Zero free parameters.}
\begin{ruledtabular}
\begin{tabular}{lccr}
Dataset & $z$ & $D_V/r_d$ (obs) & UDT dev. \\
\colrule
6dFGS  & 0.11 & 3.05  & $+4.8\%$ \\
BOSS DR12 & 0.38 & 10.27 & $+3.2\%$ \\
BOSS DR12 & 0.51 & 13.38 & $+3.9\%$ \\
BOSS DR12 & 0.61 & 15.33 & $+6.4\%$ \\
eBOSS LRG & 0.70 & 17.86 & $+2.8\%$ \\
DESI LRG1 & 0.51 & 13.62 & $+2.1\%$ \\
eBOSS QSO & 1.48 & 30.69 & $+3.7\%$ \\
DESI Ly$\alpha$ & 2.33 & 39.71 & $+1.7\%$ \\
\colrule
\multicolumn{3}{l}{RMS} & $3.8\%$ \\
\end{tabular}
\end{ruledtabular}
\end{table}

Table~\ref{tab:bao} shows the results.  The cleanest measurement
(DESI Ly$\alpha$, minimal pipeline assumptions) gives the best match
at $+1.7\%$.  The systematic positive bias in galaxy surveys is
consistent with $\Lambda$CDM fiducial contamination in data reduction.
See Fig.~\ref{fig:bao}.

\begin{figure}[t]
\includegraphics[width=\columnwidth]{figures/fig_bao_placeholder}
\caption{\label{fig:bao}%
BAO distance ratios $D_V/r_d$ versus redshift.  Points: BOSS, eBOSS,
DESI measurements.  Curve: UDT prediction with $\mug = 0.247$~Gpc$^{-1}$
and native $r_d$.}
\end{figure}

%===================================================================
\section{Reduction to General Relativity}\label{sec:gr-reduction}
%===================================================================

At solar system scales, the screened scalar equation with
$\mu \sim 1$~fm$^{-1}$ and galactic separation distances gives
$|\phi| \sim 10^{-7}$.  In this regime:
\begin{enumerate}
\item The metric reduces to the standard weak-field form
  (Sec.~\ref{sec:gr-weak}).
\item The geodesic equation gives Newtonian gravity.
\item The Shapiro time delay, light deflection, and perihelion
  precession all take their standard GR values to corrections of
  order $\phi^2 \sim 10^{-14}$, which are unmeasurable.
\item Static Coulomb electrostatics is exact (no $\phi$-corrections),
  as shown in Sec.~\ref{sec:alpha}.
\end{enumerate}
The metric~\eqref{eq:metric} is therefore indistinguishable from
standard GR at intermediate scales.  Departures appear only at
microphysical scales (where $\phi$ reaches $O(1)$ inside the cavity)
and at cosmological scales (where $\phi$ reaches $O(7)$ at the CMB
surface).

%===================================================================
\section{Angular Structure on $S^2$: The Pivot}%
\label{sec:angular}
%===================================================================

The Dirac equation on the metric~\eqref{eq:metric} separates into
radial and angular parts.  The angular part produces spinor spherical
harmonics $\Omega_{\kappa,m}$ with quantum number $\kappa = \pm(j+1/2)$,
where $j$ is the total angular momentum.  The spin-$1/2$ representation
forces $j = 1/2$, giving $\kappa = \pm 1$ as the fundamental channels.

The unit radial vector $\hat{r}$ (a metric-derived geometric operator)
couples adjacent $\kappa$ channels on $S^2$ with strength
% VERIFIED: CG §13.3
\begin{equation}\label{eq:rcoupling}
  \sum_{m,m'}\!\left|\langle\Omega_{\kappa,m}|\hat{r}|
    \Omega_{\kappa',m'}\rangle\right|^2
  = \frac{2n(n+1)}{2n+1}\,,
\end{equation}
where $n$ labels the $\kappa = -n \to \kappa = -(n+1)$ transition.
This coupling is verified numerically for $n = 1$ through 6.

\subsection{The Diophantine identity}

The pion mass coefficient $84 = \binom{9}{3}$ arises from distributing
$v = 2\ell + 1 = 3$ bosonic degrees of freedom across
$k = 2\kmax + 1 = 7$ orbit positions.  The self-consistency condition
\begin{equation}\label{eq:diophantine}
  (2j+1)^2(2\ell+1)(2\kmax+1)
  = \binom{2\ell + 2\kmax + 1}{2\ell+1}
\end{equation}
has a \emph{unique} solution for $j = 1/2$:
\begin{equation}\label{eq:jlk}
  j = \tfrac{1}{2},\quad \ell = 1,\quad \kmax = 3\,.
\end{equation}
% VERIFIED: CG §13.10
This triple determines three integers $2j+1 = 2$, $2\ell+1 = 3$,
$2\kmax + 1 = 7$, and a fourth, $2\kmax - 1 = 5$.  These four
numbers --- 2, 3, 5, 7 --- generate \emph{every} derived quantity in
the theory:
% VERIFIED: CG §17.1
\begin{table}[b]
\caption{\label{tab:multiplicities}%
The four multiplicities from the Diophantine triple
$(j, \ell, \kmax) = (1/2, 1, 3)$ and their roles.}
\begin{ruledtabular}
\begin{tabular}{lcl}
Number & Origin & Role \\
\colrule
$2 = 2j+1$ & Spin-$1/2$ & Spin bilinear, neutrino denominator \\
$3 = 2\ell+1$ & Vector ($\ell\!=\!1$) & $\pi$-power in muon, polarizations \\
$5 = 2\kmax-1$ & Closure interior & Lepton coefficient, $p$--$n$ \\
$7 = 2\kmax+1$ & Closure orbit & Cavity size, pion coefficient \\
\end{tabular}
\end{ruledtabular}
\end{table}

\subsection{The quadratic selection rule}

The product of the angular coupling~\eqref{eq:rcoupling} and the
$j$-ratio for transition $n$ is $2n(n+1)/(2n-1)$.  Setting this
equal to the universal constant $4 = (2j+1)^2$ gives a quadratic:
% VERIFIED: CG §13.4
\begin{equation}\label{eq:quadratic}
  (n-1)(n-2) = 0\,.
\end{equation}
The two roots select precisely two cross-coupled particles:
$n = 1$ (the muon) and $n = 2$ (the proton).

%===================================================================
\section{Mass Emergence}\label{sec:mass}
%===================================================================

\subsection{The Dirac operator on the metric}

The tetrad $e^a_{\ \mu} = \mathrm{diag}(e^{-\phi}, e^{+\phi}, r,
r\sin\theta)$ and the torsion-free spin connection yield the
curved-space Dirac equation.  For massless fermions ($m = 0$), the
radial system takes \emph{Form~T}:
% VERIFIED: CG §4.4
\begin{align}\label{eq:formT}
  G' + \!\left(\frac{\kappa}{r} - \phi'\right)\!G
    &= E\,e^{2\phi}\,F\,, \nonumber\\
  F' + \!\left(-\frac{\kappa}{r} - \phi'\right)\!F
    &= -E\,e^{2\phi}\,G\,,
\end{align}
where $G(r)$, $F(r)$ are the upper and lower radial spinor components,
$E$ is the eigenvalue, and $\kappa = \pm 1, \pm 2, \ldots$ labels the
angular channel.

\subsection{Boundary conditions}

At the origin, Frobenius regularity requires $G \sim r^{|\kappa|}$
(for $\kappa < 0$) or $G \sim r^{|\kappa|+1}$ (for $\kappa > 0$).
At the outer boundary $\rstar$, we impose the \emph{geometric Neumann
condition}:
% VERIFIED: CG §4.6
\begin{equation}\label{eq:neumann}
  G'(\rstar) = 0\,,
\end{equation}
which is the natural self-adjoint boundary condition of the Form-T
operator.  This condition contains \emph{zero free parameters} and
replaces earlier fitted Robin conditions.

\subsection{Locked parameters}\label{sec:params}

The three metric parameters are determined geometrically:
% VERIFIED: CG §17.2
\begin{align}
  \phiz &= -\cos(\pi/5) = -0.80902\,, \label{eq:phi0}\\
  \mu^2 &= \pi/3\,, \label{eq:mu2}\\
  \rstar &= 7 - 1/80 = 6.9875\,. \label{eq:rstar}
\end{align}
The origin of each:
\begin{itemize}
\item $\phiz$: the metric depth parameter, set by the expectation
  value $\langle G \rangle = 2/\pi$ of the upper Dirac component and
  connecting to the pentagon ($\pi/5$).
\item $\mu^2 = \pi \times \langle\cos^2\theta\rangle_{S^2} = \pi/3$:
  the scalar mass from the angular coupling on $S^2$.
\item $\rstar = (2\kmax + 1) - \mathrm{source}^2/(2\kmax - 1)
  = 7 - 1/80$: the eigenvalue-consistent cavity size, where
  $\mathrm{source} = 1/4$ is the ground-state Dirac source integral
  (Sec.~\ref{sec:alpha}).
\end{itemize}
One calibration sets the mass scale:
% VERIFIED: CG §17.4
\begin{equation}\label{eq:C}
  C = 4\pi^2 m_e \times \rstar = 140.95\;\mathrm{MeV}\,.
\end{equation}
Physical masses are then
\begin{equation}\label{eq:massformula}
  m = E_n(\kappa) \times C
\end{equation}
for all mesons except the pion (Sec.~\ref{sec:pion}).

\subsection{Ground-state eigenvalue}

At the locked parameters~\eqref{eq:phi0}--\eqref{eq:rstar}, the
$\kappa = -1$ ground-state eigenvalue is algebraic:
% VERIFIED: VR §8.15
\begin{equation}\label{eq:E1}
  E_1(\kappa\!=\!-1) = \frac{2\sqrt{2}}{3} = 0.94281\,,
\end{equation}
verified to 10~significant digits by GPU-accelerated numerical
integration.  The eigenvalue ratio
\begin{equation}\label{eq:E2E1}
  \frac{E_2}{E_1} = \frac{3\,\vgold}{I_2} \approx 5.90\,,
\end{equation}
where $\vgold = (1+\sqrt{5})/2$ and
$I_2 \equiv \intdom e^{2\phi}\,dr = 0.8230$, matches the PDG ratio
$m_\rho/m_\pi = 5.744$ to $+2.7\%$.

%===================================================================
\section{Lepton Mass Ratios}\label{sec:leptons}
%===================================================================

The angular coupling on $S^2$ (Sec.~\ref{sec:angular}) produces the
complete lepton mass formula:
% VERIFIED: CG §13.7
\begin{equation}\label{eq:lepton}
  m_n = A_n \times \pi^{Q_n - 2} \times \frac{C}{\rstar}\,,
\end{equation}
where $A_n = 5(n+1)/[2(2n+1)]$ from the rank-1 tensor coupling and
$Q_n = 2l_G + 1$ from the Gaussian measure on the complex angular
mode space.

The resulting \emph{parameter-free} mass ratios are:
% VERIFIED: CG §13.8
\begin{align}
  \frac{m_\mu}{m_e} &= \frac{20\pi^3}{3} = 206.71
    \quad(\text{PDG: }206.77,\ {-0.03\%})\,, \label{eq:mu-e}\\
  \frac{m_p}{m_e} &= 6\pi^5 = 1836.12
    \quad(\text{PDG: }1836.15,\ {-0.002\%})\,, \label{eq:p-e}\\
  \frac{m_p}{m_\mu} &= \frac{9\pi^2}{10} = 8.883
    \quad(\text{PDG: }8.880,\ {+0.03\%})\,. \label{eq:p-mu}
\end{align}
These ratios depend on \emph{nothing}: no $\phiz$, no $\mu$, no
$\rstar$, no $C$.  They are pure consequences of the angular structure
of spinor harmonics on $S^2$.

The derivation chain (Fig.~\ref{fig:leptons}) traces entirely from
the metric through the angular coupling integrals, the quadratic
selection rule~\eqref{eq:quadratic}, the Gaussian measure on the
complex mode space, and the relation $C/\rstar = 4\pi^2 m_e$.

\begin{figure}[t]
\includegraphics[width=\columnwidth]{figures/fig_leptons_placeholder}
\caption{\label{fig:leptons}%
Lepton mass ratios: $m_\mu/m_e$ and $m_p/m_e$ from the angular
coupling on $S^2$.  Horizontal bars: PDG values.  Points: UDT
predictions.  The proton--electron ratio $6\pi^5$ is accurate to
$0.002\%$.}
\end{figure}

%===================================================================
\section{The Pion}\label{sec:pion}
%===================================================================

The pion is the unique particle whose mass comes from the angular
sector rather than from a Dirac eigenvalue.  The C-self-conjugation
of the $\kappa = -1$ ground state and the bosonic occupation of the
closure orbit give:
% VERIFIED: CG §13.10
\begin{equation}\label{eq:pion}
  m_\pi = \binom{9}{3}\,\pi\,m_e = 84\pi\,m_e
  = 134.85\;\mathrm{MeV}\,,
\end{equation}
compared to PDG $135.0$~MeV ($-0.09\%$).  The coefficient
\begin{equation}
  84 = \binom{9}{3} = (2j+1)^2 \times (2\ell+1)
    \times (2\kmax + 1)
  = 4 \times 3 \times 7
\end{equation}
collects all three quantum number factors.  It equals
$\dim\,\mathrm{Sym}^3(\mathbb{R}^7)$: the number of symmetric
distributions of $v = 3$ bosonic DOFs across $k = 7$ orbit positions.

The pion also has a Dirac eigenvalue $E_1 = 2\sqrt{2}/3$
[Eq.~\eqref{eq:E1}], but the eigenvalue formula
$m_\pi = E_1 \times C = 132.9$~MeV underpredicts by $1.53\%$.  The
angular formula~\eqref{eq:pion} is $16\times$ more precise.

\begin{figure}[t]
\includegraphics[width=\columnwidth]{figures/fig_pion_placeholder}
\caption{\label{fig:pion}%
The pion mass derivation chain.  From the Diophantine triple
$(j, \ell, \kmax) = (1/2, 1, 3)$, the C-closure orbit produces
$84 = \binom{9}{3}$ configurations.  The angular formula
$m_\pi = 84\pi\,m_e$ matches PDG to $0.09\%$.}
\end{figure}

%===================================================================
\section{Hadron Spectrum}\label{sec:hadrons}
%===================================================================

\subsection{Two-domain structure}

The metric produces particle masses through eigenvalues of the Form-T
Dirac equation on two nested domains:
% VERIFIED: CG §14.1
\begin{itemize}
\item $[0, \rstar]$ --- the full cavity (mesons),
\item $[0, r_b]$ --- the sub-cavity (baryons),
\end{itemize}
where $r_b$ is the first $G$-node of the $\kappa = +1$ ground state
(the unique zero crossing of the upper spinor component).  At locked
parameters:
\begin{equation}\label{eq:goldenratio}
  \frac{\rstar}{r_b} = \vgold = \frac{1+\sqrt{5}}{2}
  \qquad(0.03\%)\,.
\end{equation}
The golden ratio partition is structural (CV $= 0.003$ across all
tested parameters).

\subsection{Meson spectrum}

% VERIFIED: VR §20.2
\begin{table}[t]
\caption{\label{tab:mesons}%
Meson predictions at locked parameters with $C = 140.95$~MeV
(electron anchor).  The pion comes from the angular sector
[Eq.~\eqref{eq:pion}]; all others from Dirac eigenvalues.}
\begin{ruledtabular}
\begin{tabular}{lcrrrr}
Meson & $\kappa, n$ & UDT (MeV) & PDG (MeV) & Error \\
\colrule
$\pi^0$       & angular & 134.85 & 135.0  & $-0.09\%$ \\
$\eta$        & $+2, 1$ & 548.0  & 547.9  & $+0.03\%$ \\
$\omega(782)$ & $-1, 2$ & 784.3  & 782.7  & $+0.20\%$ \\
$a_0(980)$    & $+1, 1$ & 973.6  & 980.0  & $-0.66\%$ \\
$\phi(1020)$  & $+3, 1$ & 1022.7 & 1019.5 & $+0.32\%$ \\
$f_2(1270)$   & $+2, 2$ & 1270.5 & 1275.5 & $-0.39\%$ \\
$a_2(1320)$   & $-1, 3$ & 1331.3 & 1318.3 & $+0.99\%$ \\
\end{tabular}
\end{ruledtabular}
\end{table}

Table~\ref{tab:mesons} shows the meson predictions.  Six mesons
(excluding the pion) match within $1\%$.  All use
$|\kappa| \leq 3$, consistent with the Diophantine triple.

\subsection{Baryon spectrum}

% VERIFIED: VR §9.1
\begin{table}[t]
\caption{\label{tab:baryons}%
Baryon predictions.  Sub-cavity eigenvalues on $[0, r_b]$ with
$C = 140.95$~MeV.  The proton is an angular-sector particle
[Eq.~\eqref{eq:p-e}], not an eigenvalue.}
\begin{ruledtabular}
\begin{tabular}{lcrrr}
Baryon & $\kappa, n$ & UDT (MeV) & PDG (MeV) & Error \\
\colrule
$p$         & angular    & 938.3  & 938.3  & $-0.002\%$ \\
$\Lambda$   & $-3, 1$    & 1119.3 & 1115.7 & $+0.32\%$ \\
$\Sigma$    & $+3, 1$    & 1174.3 & 1189.4 & $-1.27\%$ \\
$\Xi^*$     & $+4, 1$    & 1531.3 & 1531.8 & $-0.03\%$ \\
$\Omega$    & $-1, 3$    & 1681.2 & 1672.5 & $+0.52\%$ \\
$N(1680)$   & $-1, 3$    & 1681.2 & 1680.0 & $+0.07\%$ \\
\end{tabular}
\end{ruledtabular}
\end{table}

Table~\ref{tab:baryons} shows selected baryon predictions.  The
proton mass comes from the angular sector, matching PDG to $0.002\%$.
Other baryons are sub-cavity eigenvalues.  The angular correction
$m_\Omega / m_\Lambda = l_F/l_G|_{\kappa=-3} = 3/2$ is derived to
$0.06\%$.

\begin{figure}[t]
\includegraphics[width=\columnwidth]{figures/fig_hadrons_placeholder}
\caption{\label{fig:hadrons}%
Hadron mass spectrum.  Horizontal axis: PDG mass.  Vertical axis:
UDT prediction.  Diagonal line: perfect agreement.  Circles: mesons
(full cavity).  Squares: baryons (sub-cavity).  Triangle: proton
(angular sector).  The pion is marked separately.}
\end{figure}

\subsection{Three-sector summary}

The mass spectrum divides into three sectors, all from the same metric:
\begin{enumerate}
\item \textbf{Angular sector} (parameter-free):
  $e$, $\mu$, $p$, $n$, $\pi$ --- from coupling on $S^2$.
\item \textbf{Full cavity} (one calibration $C$):
  mesons --- from Dirac eigenvalues on $[0, \rstar]$.
\item \textbf{Sub-cavity} (same $C$):
  baryons --- from Dirac eigenvalues on $[0, r_b]$.
\end{enumerate}
Total: 17~particles from one metric, one ODE, one boundary condition,
one calibration.

%===================================================================
\section{The Fine-Structure Constant}\label{sec:alpha}
%===================================================================

\subsection{The identity $g^{tt}\,g^{rr} = -1/c^2$}

The metric~\eqref{eq:metric} satisfies
\begin{equation}\label{eq:gttgrr}
  g^{tt}\,g^{rr} = -\frac{e^{2\phi}}{c^2}\,e^{-2\phi}
  = -\frac{1}{c^2}\,,
\end{equation}
identically.  This reduces the static Maxwell equation to flat-space
Poisson:
\begin{equation}
  \partial_r(r^2\,\partial_r A_t) = -\rho\,.
\end{equation}
Coulomb's law $A_t = Q/(4\pi r)$ is \emph{exact} on the UDT metric
--- no $\phi$-corrections.  This is verified numerically to
$2.2 \times 10^{-16}$.

\subsection{Decomposition of $\alpha_{\mathrm{EM}}$}

The fine-structure constant decomposes into four metric-derived factors:
% VERIFIED: VR §19.4
\begin{equation}\label{eq:alpha}
  \frac{1}{\aEM}
  = \underbrace{(2j+1)^2}_{= 4}
    \times \underbrace{(2\ell+1)^2}_{= 9}
    \times \underbrace{\frac{\pi}{I_2}}_{\text{WKB gap}}
  = \frac{36\pi}{I_2} = 137.4\,,
\end{equation}
compared to $1/\alpha = 137.036$ ($+0.29\%$).  The three factors are:
\begin{itemize}
\item $(2j+1)^2 = 4$: the spin-$1/2$ bilinear has $2j+1 = 2$
  spin states; the bilinear squares this.
\item $(2\ell+1)^2 = 9$: the $\ell = 1$ vector exchange has
  $2\ell + 1 = 3$ polarizations; the exchange squares this.
\item $\pi/I_2 = 3.818$: the asymptotic eigenvalue gap, computed from
  the vacuum ODE at locked parameters.
\end{itemize}

\subsection{The bridge identity}

The product of the Dirac source integral and the ground-state
eigenvalue satisfies
\begin{equation}\label{eq:bridge}
  \mathrm{source}^2 \times E_1^2 = \frac{1}{18}
  = \frac{\ell(\ell+1)}{[(2j+1)(2\ell+1)]^2}\,,
\end{equation}
connecting the eigenvalue sector to the angular quantum numbers.  This
is the UDT analog of the Wigner--Eckart theorem.

\subsection{Proton--neutron mass splitting}

The mass difference:
% VERIFIED: VR §19.5
\begin{equation}\label{eq:pn}
  m_n - m_p = \frac{5}{4}\,\aEM\,C = 1.282\;\mathrm{MeV}
\end{equation}
compared to PDG $1.293$~MeV ($-0.87\%$).  The coefficient
$5/4 = (2\kmax - 1)/(2j+1)^2$ is the ratio of closure angular
content to spin bilinear.

%===================================================================
\section{Neutrino Mass}\label{sec:neutrino}
%===================================================================

The atmospheric neutrino mass scale follows from three metric
quantities:
% VERIFIED: VR §22.1
\begin{equation}\label{eq:mnu}
  m_\nu = \frac{\alpha^3\,m_e}{(2j+1)^2}
  = \frac{\alpha^3\,m_e}{4}
  = 0.0493\;\mathrm{eV}\,,
\end{equation}
compared to $\sqrt{\Delta m^2_{\mathrm{atm}}} = 0.0495$~eV
($-0.5\%$).  The cube $\alpha^3$ factorizes from the action into three
specific integrals~\cite{zenodo2025}:
% VERIFIED: VR §22.8
\begin{equation}\label{eq:alpha3}
  \alpha^3 = \left(\frac{1}{3\pi}\right)^{\!3}
    \times \left(\frac{1}{12}\right)^{\!3}
    \times I_2^3\,,
\end{equation}
verified to machine precision.

\subsection{Structural parallel}

The muon and neutrino share the quantum number $2\ell + 1 = 3$ but
differ in coupling:
\begin{align}
  m_\mu &= \frac{20}{3}\,\pi^3\,m_e
    &&(\text{angular coupling }\pi)\,,\\
  m_\nu &= \frac{1}{4}\,\alpha^3\,m_e
    &&(\text{EM coupling }\alpha)\,.
\end{align}
The nine-order-of-magnitude hierarchy $m_\mu/m_\nu \approx 2 \times
10^9$ is combinatorial, arising from $(\pi/\alpha)^3$.

\subsection{The Hodge dual}

The neutrino mass formula lives in the exterior algebra of the
9-dimensional angular base space
($\dim = 2(\ell + \kmax) + 1 = 9$).  The pion occupies grade 3
($\Lambda^3$, symmetric, $\binom{9}{3} = 84$ configurations), while
the neutrino occupies grade 6 ($\Lambda^6 \cong {*}\Lambda^3$,
antisymmetric, accessed through three 2-form EM couplings):
\begin{equation}
  \Lambda^2 \wedge \Lambda^2 \wedge \Lambda^2 \to \Lambda^6
  \cong {*}\Lambda^3\,.
\end{equation}
The pion (heavy, grade~3) and the neutrino (light, grade~6) are Hodge
duals on the same base space.

\begin{figure}[t]
\includegraphics[width=\columnwidth]{figures/fig_neutrino_placeholder}
\caption{\label{fig:neutrino}%
Neutrino mass derivation.  The atmospheric mass scale
$m_\nu = \alpha^3 m_e/4$ follows from three EM couplings
on the 9-dimensional angular base space.}
\end{figure}

%===================================================================
\section{PMNS Mixing Angles}\label{sec:pmns}
%===================================================================

The three neutrino mixing angles are determined by the same quantum
numbers $(2j+1, 2\ell+1, 2\kmax+1) = (2, 3, 7)$:
% VERIFIED: VR §22.7
\begin{align}
  \sin^2\theta_{12} &= \frac{(2j+1)^2}{(2j+1)^2 + (2\ell+1)^2}
    = \frac{4}{13}\,, \label{eq:theta12}\\
  \sin^2\theta_{23} &= \frac{(2j+1)^2}{2\kmax + 1}
    = \frac{4}{7}\,, \label{eq:theta23}\\
  \sin^2\theta_{13} &= \frac{1}{(2\ell+1)^2(2\kmax - 1)}
    = \frac{1}{45}\,. \label{eq:theta13}
\end{align}

\begin{table}[t]
\caption{\label{tab:pmns}%
PMNS mixing angles.  UDT predictions from quantum number fractions
vs.\ NuFIT~5.2 (normal ordering).  All within $1\sigma$.}
\begin{ruledtabular}
\begin{tabular}{lcccr}
Angle & UDT & Expt.\ & $\sigma$ & Dev. \\
\colrule
$\sin^2\theta_{12}$ & $4/13 = 0.3077$
  & $0.304 \pm 0.012$ & 0.31 & $+1.2\%$ \\
$\sin^2\theta_{23}$ & $4/7 = 0.5714$
  & $0.573 \pm 0.018$ & 0.09 & $-0.3\%$ \\
$\sin^2\theta_{13}$ & $1/45 = 0.02222$
  & $0.02220 \pm 0.00062$ & 0.04 & $+0.1\%$ \\
\end{tabular}
\end{ruledtabular}
\end{table}

The denominators have physical meaning:
\begin{itemize}
\item $13 = 4 + 9 = (2j+1)^2 + (2\ell+1)^2$: spin-orbital partition.
\item $7 = 2\kmax + 1$: closure orbit size.
\item $45 = 9 \times 5 = (2\ell+1)^2 \times (2\kmax - 1)$: orbital
  barrier times angular content.
\end{itemize}
The PMNS matrix constructed from these angles is exactly unitary
(verified numerically).

\begin{figure}[t]
\includegraphics[width=\columnwidth]{figures/fig_pmns1_placeholder}
\caption{\label{fig:pmns1}%
PMNS mixing angles: UDT predictions (vertical lines) vs.\ NuFIT
5.2 distributions.  All three predictions fall within the $1\sigma$
bands.}
\end{figure}

\begin{figure}[t]
\includegraphics[width=\columnwidth]{figures/fig_pmns2_placeholder}
\caption{\label{fig:pmns2}%
The quantum number denominators (13, 7, 45) in the PMNS formulas and
their decomposition into the four multiplicities $\{2, 3, 5, 7\}$.}
\end{figure}

%===================================================================
\section{The Generalized Bridge Identity}\label{sec:bridge}
%===================================================================

The bridge identity~\eqref{eq:bridge} is the first row of a
three-row structure that transitions between the spin ($j$) and
orbital ($\ell$) sectors:
% VERIFIED: VR §22.2
\begin{equation}\label{eq:gbridge}
  \mathrm{source}(\kappa\!=\!-n)^2 \times E_1(\kappa\!=\!-n)^2
  = B(n)\,.
\end{equation}

\begin{table}[t]
\caption{\label{tab:bridge}%
Generalized bridge identity: the $j$--$\pi$--$\ell$ transition across
$|\kappa| = 1, 2, 3$.}
\begin{ruledtabular}
\begin{tabular}{lccccr}
$|\kappa|$ & source$\cdot E_1$ & source$^2\!\cdot E_1^2$ &
  Algebraic & Sector & Match \\
\colrule
1 & $\sqrt{2}/6$ & $1/18$ & $1/[2(2\ell+1)^2]$ & orbital & 0.21\% \\
2 & $\pi/6$ & $\pi^2/36$ & $(\pi/[(2j+1)(2\ell+1)])^2$
  & mixed & 0.11\% \\
3 & $\sqrt{2}/4$ & $1/8$ & $1/[2(2j+1)^2]$ & spin & 0.045\% \\
\end{tabular}
\end{ruledtabular}
\end{table}

Table~\ref{tab:bridge} shows the transition: at $|\kappa| = 1$,
the orbital quantum number $\ell$ governs; at
$|\kappa| = 3 = \kmax$, the spin quantum number $j$ takes over; at
$|\kappa| = 2$, both participate and $\pi$ enters as the $S^2$
circumference factor.

\begin{figure}[t]
\includegraphics[width=\columnwidth]{figures/fig_bridge_placeholder}
\caption{\label{fig:bridge}%
The generalized bridge identity.  The product
$\mathrm{source}^2 \times E_1^2$ transitions from the orbital sector
($\ell$-dominated) at $|\kappa| = 1$ through the mixed sector at
$|\kappa| = 2$ to the spin sector ($j$-dominated) at $|\kappa| = 3$.}
\end{figure}

%===================================================================
\section{CMB Peak Positions}\label{sec:cmb}
%===================================================================

\subsection{WKB peak law on the sourced profile}

The scalar wave equation on the cosmological $\phi$ profile yields
CMB peak positions via the WKB formula:
% VERIFIED: VR §2.2
\begin{equation}\label{eq:wkb}
  \ell_n = \ell_A\!\left(\nu_n + \frac{q_1}{\nu_n}
    + \frac{q_2}{\nu_n^3}\right)\,,
  \qquad \nu_n = n + c_{\mathrm{eff}}\,,
\end{equation}
where $\ell_A = 220/F_1$ is the acoustic scale, $c_{\mathrm{eff}} =
\arctan(\eta)/\pi$ is the phase offset, and $q_1$, $q_2$ are curvature
corrections from $\phi$ derivatives at $\rstar$.  All quantities are
derived from the metric.

\subsection{Results}

% VERIFIED: VR §2.3
\begin{table}[t]
\caption{\label{tab:cmb}%
CMB peak positions.  Best fit with $\mug = 0.248$~Gpc$^{-1}$
(1~free parameter).  No $H_0$, $\Omega_m$, $\Omega_\Lambda$,
$\Omega_b$, $n_s$, or $\tau$.}
\begin{ruledtabular}
\begin{tabular}{lcrc}
Peak & $\ell_{\mathrm{pred}}$ & $\ell_{\mathrm{obs}}$ (Planck)
  & Error \\
\colrule
1 & 220.0 & 220.0 & 0.00\% (anchored) \\
2 & 536.8 & 537.5 & $-0.13\%$ \\
3 & 831.1 & 810.8 & $+2.50\%$ \\
4 & 1133.3 & 1120.9 & $+1.10\%$ \\
5 & 1438.8 & 1444.2 & $-0.38\%$ \\
6 & 1745.9 & 1776.0 & $-1.69\%$ \\
7 & 2054.0 & 2081.0 & $-1.30\%$ \\
\colrule
\multicolumn{3}{l}{RMS (peaks 1--7)} & 1.32\% \\
\end{tabular}
\end{ruledtabular}
\end{table}

With one free parameter ($\mug$), the RMS error over seven peaks is
$1.32\%$ (Table~\ref{tab:cmb}).  The Planck data processing pipeline
uses $\Lambda$CDM transfer functions for foreground subtraction and
beam deconvolution, introducing an estimated 2--3\% model
contamination.  A model-independent CMB reduction would be needed for
a fairer comparison.

\subsection{Cross-sector link}

The acoustic scale $\ell_A$ is algebraically related to microphysical
eigenvalues:
\begin{equation}
  \ell_A = \frac{2\pi\rstar(E_2/E_1)}{I_2}
  = \frac{6\pi\rstar\vgold}{I_2^2}\,,
\end{equation}
connecting the CMB peak spacing to the meson mass ratio.

\begin{figure}[t]
\includegraphics[width=\columnwidth]{figures/fig_cmb_placeholder}
\caption{\label{fig:cmb}%
CMB angular power spectrum: peak positions.  Vertical lines: Planck
2018 measured positions.  Points: UDT predictions from WKB on the
sourced $\phi$ profile.  RMS: 1.32\%.}
\end{figure}

%===================================================================
\section{JWST Predictions}\label{sec:jwst}
%===================================================================

In UDT, the universe is static and eternal.  There is no Big Bang, no
``first epoch'' of galaxy formation, and no age--redshift
relationship.  A galaxy at $z = 12$ is not younger than one at
$z = 2$ --- it is at a different radial position in the $\phi$
profile, experiencing different dilation.

This produces a specific prediction:
\begin{quote}
\emph{Massive, structurally mature galaxies should exist at all
redshifts, with no systematic trend toward ``youth'' at high~$z$.}
\end{quote}
The JWST observations of massive galaxies at $z > 10$
--- widely described as ``impossibly early'' in the $\Lambda$CDM
framework~\cite{Labbe2023,Boylan2023} --- are precisely what UDT
predicts.  The ``impossible early galaxies'' are simply galaxies seen
through a deeper part of the $\phi$ gradient.  In UDT, there is no
paradox.

This prediction was inherent in the metric from the start and
predates the JWST data.  It is qualitative (the galaxy population
statistics have not been quantitatively modeled), but the qualitative
content --- mature galaxies at all redshifts --- is confirmed by JWST
observations and is in direct tension with $\Lambda$CDM expectations.

%===================================================================
\section{Rotation Curves}\label{sec:rotcurves}
%===================================================================

\subsection{First-principles derivation}

The geodesic equation on the metric~\eqref{eq:metric} gives the
circular velocity:
% VERIFIED: VR §5.1
\begin{equation}\label{eq:vcirc}
  v^2(r) = c^2\,r\,|\phi'(r)|\,e^{-2\phi(r)}\,.
\end{equation}
At galactic scales, $\mu_g r \sim 10^{-5}$, so the screening term
$\mu_g^2\phi$ is suppressed by $O(\mu_g^2 r^2) \sim 10^{-10}$.  In
the weak-field limit:
\begin{equation}\label{eq:newt}
  v^2_{\mathrm{UDT}}(r) = \frac{G\,M_{\mathrm{bar}}(<r)}{r}
  = v^2_{\mathrm{Newton}}(r)\,.
\end{equation}
UDT recovers Newtonian gravity exactly at galactic scales.  The flat
rotation curve problem is open in UDT, as in standard GR.

\subsection{Ancient baryonic remnants}

In UDT's static, eternal universe, galaxies accumulate dead stellar
remnants (white dwarfs, neutron stars, black holes) over
100--1000~Gyr.  Supernova kicks naturally form halos.  A
one-parameter ($T_{\mathrm{age}}$) model applied to 175~SPARC
galaxies achieves 96\%~improvement over baryons-only for massive
galaxies, within $1.76\times$ of NFW dark matter halos with half the
free parameters~\cite{zenodo2025}.

This mechanism requires no new fields, no modified geodesics, and no
dark matter --- only standard stellar evolution on longer timescales.

\begin{figure}[t]
\includegraphics[width=\columnwidth]{figures/fig_sparc_placeholder}
\caption{\label{fig:sparc}%
SPARC rotation curve fits.  Left: UDT baryons-only.  Right: UDT with
ancient baryonic remnants ($T_{\mathrm{age}}$ model).  The remnant
mechanism closes most of the gap between baryons-only and NFW fits.}
\end{figure}

%===================================================================
\section{Black Holes}\label{sec:blackholes}
%===================================================================

\subsection{The positive-$\phiz$ branch}

The vacuum classification (Sec.~\ref{sec:vacuum}) produces black
holes from $\phiz > 0$.  In this branch, $\phi(r) \to +\infty$ at
a finite radius $r_c$, where
% VERIFIED: CG §15.2
\begin{equation}
  g_{tt} \to 0,\qquad g_{rr} \to +\infty\,,
\end{equation}
and the Misner--Sharp mass gives $r_c = 2\,\mMS$ --- the Schwarzschild
radius.  UDT and GR agree on the horizon location.

\subsection{Differences from Schwarzschild}

\begin{table}[b]
\caption{\label{tab:bh}%
Structural differences between UDT and Schwarzschild black holes.}
\begin{ruledtabular}
\begin{tabular}{lcc}
Property & Schwarzschild & UDT \\
\colrule
Horizon & $g_{tt} = 0$ at $r = 2M$ & $g_{tt} \to 0$ asymptotically \\
Interior & Sign swap & No sign swap \\
Singularity & At $r = 0$ & At $r = r_c$ (boundary) \\
Center & Singular & Regular ($\phi_0$ finite) \\
\end{tabular}
\end{ruledtabular}
\end{table}

The UDT black hole (Table~\ref{tab:bh}) has no interior singularity.
The center is regular ($\phi(0) = \phiz$), and the ``singularity'' is
at the boundary surface $r_c$.  There is no causal disconnection.

\subsection{Matter dissolution and recycling}

A particle cavity ($\phiz < 0$) approaching a black hole
($\phiz > 0$) experiences a locally increasing $\phi$.  When the
effective $\phiz$ crosses zero, the eigenvalue spectrum ceases to
exist (Sec.~\ref{sec:mass}): the particle dissolves.  This occurs
\emph{outside} the horizon, providing a natural mechanism for matter
recycling in a static universe.

%===================================================================
\section{Predictions and Falsification}\label{sec:falsify}
%===================================================================

A theory with 30+ predictions must specify the conditions under which
it would be falsified.  Table~\ref{tab:falsify} lists concrete
falsification criteria with named experiments.

% VERIFIED: project design
\begin{table*}[t]
\caption{\label{tab:falsify}%
Falsification criteria.  Each row identifies a prediction, the
threshold for falsification, and the experiment or measurement that
would test it.}
\begin{ruledtabular}
\begin{tabular}{llll}
Prediction & Falsification criterion & Experiment/Measurement \\
\colrule
$m_p/m_e = 6\pi^5$ & Ratio deviates $> 0.01\%$ from $6\pi^5$
  & CODATA precision mass measurements \\
$m_\mu/m_e = 20\pi^3/3$ & Ratio deviates $> 0.1\%$
  & Muon $g{-}2$ experiments (Fermilab, J-PARC) \\
$m_\pi = 84\pi\,m_e$ & Ratio deviates $> 0.5\%$
  & PIENU, precision pion mass \\
$1/\aEM = 36\pi/I_2$ & Deviation $> 1\%$
  & Electron $g{-}2$ (Harvard) \\
$\sin^2\theta_{12} = 4/13$ & Deviation $> 3\sigma$
  & JUNO, DUNE \\
$\sin^2\theta_{23} = 4/7$ & Deviation $> 3\sigma$
  & Hyper-Kamiokande, DUNE \\
$\sin^2\theta_{13} = 1/45$ & Deviation $> 3\sigma$
  & JUNO \\
$m_\nu = \alpha^3 m_e/4$ & Direct mass $> 0.06$~eV
  or $< 0.04$~eV & KATRIN, Project~8 \\
$m_n - m_p = (5/4)\aEM C$ & Deviation $> 2\%$
  & Lattice QCD independent verification \\
Static universe & Detection of cosmological expansion
  at $> 5\sigma$ & Sandage--Loeb test (ELT) \\
Mature galaxies at all $z$ & Systematic absence of
  mature galaxies at $z > 15$ & JWST, Roman Space Telescope \\
No dark matter particle & Detection of WIMP or axion
  & LZ, XENONnT, ADMX \\
$c^2 = 2GM/\rstar$ & Machian relation fails at $> 1\%$
  & Improved BAO + CMB joint analysis \\
\end{tabular}
\end{ruledtabular}
\end{table*}

The most immediate tests:
\begin{enumerate}
\item \textbf{JUNO} (2025--2030): will measure
  $\sin^2\theta_{12}$ to $\pm 0.005$, testing the prediction
  $4/13 = 0.3077$ at $> 3\sigma$ sensitivity.
\item \textbf{KATRIN/Project~8}: direct neutrino mass measurements
  approaching the $0.04$--$0.06$~eV range where
  $m_\nu = \alpha^3 m_e/4 = 0.049$~eV lies.
\item \textbf{Sandage--Loeb test} (ELT drift spectroscopy): will
  directly test cosmic expansion vs.\ static geometry over 10--20~year
  baselines.
\item \textbf{JWST/Roman}: continued high-$z$ surveys will either
  find or rule out massive mature galaxies at $z > 15$.
\end{enumerate}

\begin{figure}[t]
\includegraphics[width=\columnwidth]{figures/fig_falsify_placeholder}
\caption{\label{fig:falsify}%
Falsification landscape.  Each point represents a prediction; the
horizontal axis shows current experimental precision; the vertical
axis shows the UDT prediction.  Arrows indicate experiments that
will improve precision into the falsification range within 5~years.}
\end{figure}

%===================================================================
\section{Discussion}\label{sec:discussion}
%===================================================================

\subsection{Summary of the framework}

The results presented in this paper flow from a single line element
[Eq.~\eqref{eq:metric}] through a chain of standard mathematical
operations: tetrad construction, spin connection, covariant Dirac
equation, Sturm--Liouville eigenvalue theory, angular coupling
integrals on $S^2$, and the WKB approximation for the scalar wave
equation.  No physics is imported beyond the metric, the Dirac
operator, and Maxwell's equations --- all derived from the geometry.

The three parameters $(\phiz, \mu^2, \rstar)$ are determined by
angular quantum numbers on $S^2$, and the single calibration $C$
sets the mass scale via the electron mass.  From this minimal input,
the theory produces:
\begin{itemize}
\item Parameter-free lepton mass ratios accurate to $0.002$--$0.03\%$.
\item The pion mass from the Diophantine identity $84 = \binom{9}{3}$
  at $0.09\%$.
\item A meson--baryon spectrum matching 17~particles within 1\%.
\item The fine-structure constant at $0.29\%$.
\item Three PMNS mixing angles, all within $1\sigma$.
\item The atmospheric neutrino mass at $0.5\%$.
\item SNe distances at $0.166$~mag RMS.
\item CMB peak positions at $1.32\%$ RMS.
\item BAO distances at $3.8\%$ RMS.
\end{itemize}

\subsection{The static universe}

The metric~\eqref{eq:metric} is static: there is no time dependence,
no expansion, no Big Bang.  Cosmological redshift arises from the
spatial gradient of $\phi$ [Eq.~\eqref{eq:redshift}], not from the
expansion of space.  The CMB is the last-scattering surface at the
radial distance where $\phi(\rstar) = \ln(1+z_{\mathrm{CMB}}) = 7.004$.

This is a legitimate general-relativistic solution.  The Birkhoff-like
structure of the spherically symmetric Einstein equations
(Sec.~\ref{sec:metric}) admits both static and dynamic solutions; we
work with the static branch.  Whether the universe is expanding or
static is an observational question, and we note that the Sandage--Loeb
test~\cite{Sandage1962,Loeb1998} --- a model-independent measurement
of cosmic drift --- has not yet been performed at the required
sensitivity.

\subsection{The role of $\phi$}

The scalar field $\phi(r)$ plays a dual role:
\begin{enumerate}
\item \textbf{Microphysical}: the dilation field creates a potential
  well that confines Dirac eigenstates, producing the particle mass
  spectrum.
\item \textbf{Cosmological}: the same field, integrated over
  cosmological distances, produces the redshift-distance relation,
  CMB peak positions, and BAO structure.
\end{enumerate}
The two regimes are cleanly separated: the microphysical cavity has
$\rstar \approx 7$ in natural units ($\sim 7$~fm), while the
cosmological boundary has $\rstar \approx 10$~Gpc.  The same equation
[Eq.~\eqref{eq:kg}] governs both, with different parameters
($\mu$ vs.\ $\mug$) and different sources.

\subsection{What this theory does not explain}

We note the following open problems:
\begin{itemize}
\item \textbf{Rotation curves}: UDT recovers Newtonian gravity at
  galactic scales.  The flat rotation curve problem is open, though
  the ancient baryonic remnant mechanism is promising.
\item \textbf{CMB peak heights}: the full power spectrum (including
  peak amplitudes) is not yet competitive with $\Lambda$CDM.
\item \textbf{The tau lepton}: no mass formula has been derived.
\item \textbf{The CP-violating phase}: the fourth PMNS parameter
  $\delta$ is not yet predicted.
\item \textbf{Cross-sector closure}: the mechanism connecting
  microphysical and cosmological parameters is identified (shared
  geometric constants) but not fully derived from a single action.
\end{itemize}

\subsection{Comparison with $\Lambda$CDM}

The standard cosmological model uses six parameters
($H_0$, $\Omega_b$, $\Omega_m$, $n_s$, $A_s$, $\tau$) to fit
cosmological data.  UDT uses one cosmological parameter ($\mug$) and
achieves comparable results on peak positions, while falling short on
peak heights and the full CMB power spectrum.

On particle physics, $\Lambda$CDM has no predictions (particle masses
are free parameters in the Standard Model).  UDT derives 17+~particle
masses and three mixing angles from the metric --- content that the
Standard Model takes as input.

The comparison is not symmetric: UDT attempts to unify particle
physics and cosmology from a single geometric object, while
$\Lambda$CDM addresses only cosmology and treats particle physics
as external input.

%===================================================================
\section{Conclusion}\label{sec:conclusion}
%===================================================================

We have shown that the static, spherically symmetric metric
\begin{equation*}
  ds^2 = -e^{-2\phi(r)}\,c^2\,dt^2 + e^{2\phi(r)}\,dr^2
         + r^2\,d\Omega^2
\end{equation*}
with a single screened scalar field $\phi(r)$ produces more than
30~quantitative predictions spanning particle masses, coupling
constants, neutrino parameters, and cosmological observables.  All
predictions flow from the metric through standard mathematical
operations, with three geometric parameters and one mass anchor.

The key results --- $m_p/m_e = 6\pi^5$ at $0.002\%$,
$m_\pi = 84\pi\,m_e$ at $0.09\%$, $1/\alpha = 36\pi/I_2$ at
$0.29\%$, three PMNS angles within $1\sigma$, and the atmospheric
neutrino mass at $0.5\%$ --- are either correct or very precisely
wrong.  The falsification criteria (Table~\ref{tab:falsify}) specify
the experiments that will distinguish these two possibilities within
the next decade.

This is a modest extension of general relativity: one metric, one
scalar field, one principle (everything from the geometry).  The
metric was published in September 2025.  All numerical codes,
derivations, and datasets are available at \cite{zenodo2025}.

\begin{acknowledgments}
The authors acknowledge the use of the Pantheon+ supernova
compilation~\cite{Pantheon2022}, Planck 2018 CMB
data~\cite{Planck2018}, BOSS and DESI BAO
measurements~\cite{DESI2024}, SPARC rotation curve
data~\cite{SPARC2016}, the Particle Data Group
compilations~\cite{PDG2024}, and the NuFIT neutrino oscillation
global fit~\cite{NuFIT}.  Numerical computations were performed on a
Tesla V100 GPU using PyTorch for vectorized ODE integration.  This
work was supported by no external funding.  All code is available
under CC~BY-NC-ND~4.0 at \cite{zenodo2025}.
\end{acknowledgments}

\appendix

%===================================================================
\section{The Einstein Tensor}\label{app:einstein}
%===================================================================

For the metric~\eqref{eq:metric-full}, the nonzero mixed components
of the Einstein tensor are:
\begin{align}
  G^t_{\ t} &= -\frac{1}{r^2}\left[1 - e^{-2\phi}(1 - 2r\phi')
    \right]\,, \label{eq:Gtt}\\
  G^\theta_{\ \theta} &= e^{-2\phi}\left[-\phi'' + 2(\phi')^2
    - \frac{2\phi'}{r}\right]\,, \label{eq:Gthth}
\end{align}
with $G^r_{\ r} = G^t_{\ t}$ (identically) and
$G^\varphi_{\ \varphi} = G^\theta_{\ \theta}$.

The Bianchi identity $\nabla_\mu G^{\mu\nu} = 0$ gives the integral
form
\begin{equation}
  G^t_{\ t} = \frac{1}{r^2}\left[C_0 + I(r)\right]\,,
\end{equation}
where $C_0 = e^{-2\phiz} - 1$ and
$I(r) = -2\mu^2\int_0^r s\,\phi(s)\,ds$.  Both decompositions are
verified numerically to residual $< 2 \times 10^{-9}$.

%===================================================================
\section{Misner--Sharp Mass and the Machian Relation}%
\label{app:misner}
%===================================================================

The Misner--Sharp quasi-local mass for the metric is
\begin{equation}\label{eq:ms}
  m(r) = \frac{c^2\,r}{2G}\bigl(1 - e^{-2\phi(r)}\bigr)\,.
\end{equation}
At the boundary $r = \rstar$ with total mass $M$:
\begin{equation}
  c^2 = \frac{2GM}{\rstar(1 - e^{-2\phi(\rstar)})}\,.
\end{equation}
In the deep-well limit $|\phi(\rstar)| \gg 1$
($e^{-2\phi_*} \to 0$):
\begin{equation}\label{eq:machian}
  c^2 \approx \frac{2GM}{\rstar}\,,
\end{equation}
verified to $< 0.01\%$ at BAO-constrained parameters.  The cosmological
boundary value
$\phi(\rstar) = \ln(1 + z_{\mathrm{CMB}}) = 7.004$ is algebraically
fixed by the combination of the redshift definition~\eqref{eq:redshift}
and the Misner--Sharp relation~\eqref{eq:ms}.

%===================================================================
\section{Dirac Stress-Energy Tensor}\label{app:tmunu}
%===================================================================

For a massless Dirac eigenstate $(G, F, E, \kappa)$ on the UDT metric,
the stress-energy components (after angular integration) are:
% VERIFIED: CG §9.5
\begin{align}
  T^t_{\ t}(r) &= -\frac{E\,e^{\phi}(G^2 + F^2)}{4\pi\,r^2}\,,
    \label{eq:Ttt}\\
  T^r_{\ r}(r) &= T^t_{\ t}
    + \frac{2\kappa\,e^{-\phi}\,G\,F}{4\pi\,r^3}\,,
    \label{eq:Trr}\\
  T^\theta_{\ \theta}(r) &= -\frac{1}{2}(T^t_{\ t} + T^r_{\ r})\,.
    \label{eq:Tthth}
\end{align}
The trace vanishes exactly: $T^\mu_{\ \mu} = 0$ (for $m = 0$).  The
ratio $R \equiv T^r_{\ r}/T^t_{\ t}$ ranges from 0.094 (ground state)
to 0.841 ($n = 3$ excited), approaching 1 in the ultrarelativistic
limit.  The constraint $G^t_{\ t} = G^r_{\ r}$ requires $R = 1$ for
total matter, which is achieved by including the electromagnetic
sector ($R_{\mathrm{EM}} = 1$ exactly).

%===================================================================
\section{The Vacuum Classification Theorem}\label{app:vacuum}
%===================================================================

Define $y(r) := e^{-2\phi(r)} > 0$.  The vacuum flux system is
equivalent to
\begin{equation}\label{eq:yode}
  y'' + \frac{2}{r}\,y' = \mu^2\ln y\,,\quad
  y(0) = e^{-2\phiz},\quad y'(0) = 0\,.
\end{equation}

\textbf{Theorem.} For $\mu > 0$ and regular origin data:
\begin{enumerate}
\item[(i)] $\phiz > 0$: $y$ reaches zero at finite $r_c$
  ($\phi \to +\infty$, black hole horizon).
\item[(ii)] $\phiz = 0$: $y \equiv 1$, $\phi \equiv 0$ (trivial).
\item[(iii)] $\phiz < 0$: $y$ is globally defined, strictly
  increasing, $y \to \infty$ ($\phi \to -\infty$).
\end{enumerate}

\textbf{Proof.}  Case~(i): $y_0 < 1$, so $\ln y < 0$, driving $y$
downward.  The quadratic bound
$y(r) \leq y_0 - (\mu^2 a/6)\,r^2$ with $a = -\ln y_0 > 0$
guarantees $y = 0$ at finite $r_c$.  Case~(iii): $y_0 > 1$, so
$\ln y > 0$, driving $y$ upward.  Global existence follows from
$y(r) \leq y_0\exp(\mu^2 r^2/6)$.  A finite limit $L > 1$ leads to
$y' \sim cr$ for large $r$, contradicting convergence.  $\square$

%===================================================================
\section{The Separatrix Theorem}\label{app:separatrix}
%===================================================================

For the sourced scalar equation with $\varepsilon$-decomposition
$S(r) = \mu^2\phi(1 - \varepsilon)$, the \emph{null-source curve}
satisfies $J' = 0$ everywhere.  Setting $u = \phi'(r)$ gives a
Bernoulli ODE with solution:
\begin{equation}\label{eq:separatrix}
  \phi'_{\mathrm{sep}}(r) = \frac{1}{r(2 + Cr)}\,,
\end{equation}
where $C$ is an integration constant (distinct from the mass
calibration).  Any profile with $\phi'$ above this curve is in the
drive regime; below is in the brake regime.  The sign-change radius
$r_{sc}$ (maximum of the scalar flux) satisfies
$J'(r_{sc}) = 0$ and is determined entirely by the profile
coefficients, with no external input.

%===================================================================
% BIBLIOGRAPHY
%===================================================================

\begin{thebibliography}{99}

\bibitem{zenodo2025}
C.~Rotter and A.~Watts,
``Universal Dilation Theory: metric, operators, and predictions,''
Zenodo (2025).
\url{https://doi.org/10.5281/zenodo.XXXXXXX}

\bibitem{PDG2024}
R.~L.~Workman \textit{et al.} (Particle Data Group),
Prog.\ Theor.\ Exp.\ Phys.\ \textbf{2022}, 083C01 (2022);
updated 2024.

\bibitem{Planck2018}
N.~Aghanim \textit{et al.} (Planck Collaboration),
Astron.\ Astrophys.\ \textbf{641}, A6 (2020).

\bibitem{Pantheon2022}
D.~Scolnic \textit{et al.},
Astrophys.\ J.\ \textbf{938}, 113 (2022).

\bibitem{DESI2024}
DESI Collaboration,
arXiv:2404.03002 (2024).

\bibitem{NuFIT}
I.~Esteban \textit{et al.},
J.\ High Energy Phys.\ \textbf{2020}, 178 (2020);
NuFIT~5.2 (2022), \url{http://www.nu-fit.org}.

\bibitem{SPARC2016}
F.~Lelli, S.~S.~McGaugh, and J.~M.~Schombert,
Astron.\ J.\ \textbf{152}, 157 (2016).

\bibitem{Will2014}
C.~M.~Will,
Living Rev.\ Relativity \textbf{17}, 4 (2014).

\bibitem{Sandage1962}
A.~Sandage,
Astrophys.\ J.\ \textbf{136}, 319 (1962).

\bibitem{Loeb1998}
A.~Loeb,
Astrophys.\ J.\ \textbf{499}, L111 (1998).

\bibitem{Labbe2023}
I.~Labb\'{e} \textit{et al.},
Nature \textbf{616}, 266 (2023).

\bibitem{Boylan2023}
M.~Boylan-Kolchin,
Nat.\ Astron.\ \textbf{7}, 731 (2023).

\end{thebibliography}

\end{document}
